% =================================================================
% EV Yük Dengeleme — Birleşik Teknik & İş Planı (v1)
% 3 Kurucu, 12 Ay Hedef
%
% Derleme: pdflatex teknik_is_plani.tex (2 kez — toc için)
% =================================================================
\documentclass[11pt,a4paper]{article}

\usepackage[T1]{fontenc}
\usepackage[utf8]{inputenc}
\usepackage[turkish]{babel}
\usepackage[a4paper, margin=2.2cm]{geometry}
\usepackage{titlesec}
\usepackage{enumitem}
\usepackage{xcolor}
\usepackage{tabularx}
\usepackage{booktabs}
\usepackage{longtable}
\usepackage{array}
\usepackage{hyperref}
\usepackage{parskip}
\usepackage{tcolorbox}
\usepackage{listings}

\hypersetup{
    colorlinks=true,
    linkcolor=blue!50!black,
    urlcolor=blue!60!black
}

\titleformat{\section}{\Large\bfseries\color{blue!40!black}}{\thesection}{0.6em}{}
\titleformat{\subsection}{\large\bfseries\color{blue!30!black}}{\thesubsection}{0.5em}{}
\titleformat{\subsubsection}{\normalsize\bfseries\color{blue!20!black}}{\thesubsubsection}{0.4em}{}

\setlist[itemize]{leftmargin=1.2em, itemsep=2pt, topsep=3pt}
\setlist[enumerate]{leftmargin=1.4em, itemsep=2pt, topsep=3pt}

\newtcolorbox{notebox}[1][]{colback=yellow!8, colframe=yellow!50!black, boxrule=0.4pt, arc=2pt, left=6pt, right=6pt, top=3pt, bottom=3pt, #1}
\newtcolorbox{redbox}[1][]{colback=red!5, colframe=red!50!black, boxrule=0.4pt, arc=2pt, left=6pt, right=6pt, top=3pt, bottom=3pt, #1}
\newtcolorbox{greenbox}[1][]{colback=green!5, colframe=green!50!black, boxrule=0.4pt, arc=2pt, left=6pt, right=6pt, top=3pt, bottom=3pt, #1}

\title{\textbf{EV Yük Dengeleme — Teknik \& İş Planı}\\
       \large Birleşik Yol Haritası — 3 Kurucu, 12 Ay\\
       \normalsize v1, \today}
\author{}
\date{}

\begin{document}
\maketitle
\tableofcontents
\newpage

% ===================================================================
\section{Yönetici Özeti}

Bu doküman, \texttt{simulation\_v3.py} prototipinden yola çıkarak Türkiye'de ticari ürün haline getirme planının \textbf{hem teknik hem de iş} cephesini birleştirir. Hedef: 12 ay sonunda en az 5 ödeyen kurumsal müşteri, çalışan bulut platformu, ölçülebilir ROI raporu, Seed turuna hazır metrikler.

\begin{notebox}
\textbf{Şu anda neredeyiz (TRL 3 — Lab Prototipi):}
\begin{itemize}
    \item \texttt{simulation\_v3.py}: 5 algoritmalı (Algoritmasız / Yönetimli / SRPT / Water-Filling / Dinamik Adil), 5 saha senaryolu (AVM, ofis, otel, hastane, havalimanı) çalışan dakika çözünürlüklü simülatör.
    \item Sentetik veri (\texttt{dataset.json}) ile test edilmiş; gerçek charger, gerçek tarife, gerçek müşteri yok.
    \item Tek-süreçli (mono script), tek-saha, tek-gün simülasyon.
    \item Excel + matplotlib raporları üretiliyor.
\end{itemize}
\textbf{Eksikler — 12 ayda kapatılacak:} Maliyet modeli, OCPP entegrasyonu, gerçek-zamanlı veri toplama, yük tahmini, çoklu kiracı bulut mimarisi, güvenlik/uyum, harmonik analiz (Faz 4).
\end{notebox}

\textbf{Stratejik öneri (TL;DR):}
\begin{itemize}
    \item \textbf{Salt yazılım} (donanım üretmiyoruz). OCPP üzerinden mevcut chargerlara bağlanıyoruz.
    \item \textbf{Birincil müşteri = saha sahibi} (AVM, otopark, lojistik depo, otel). CPO ikinci faz, devlet üçüncü faz.
    \item \textbf{İlk para = TÜBİTAK 1512 BiGG} (200K TL hibe, geri ödemesiz). Pre-seed Ay 5--8 arası.
    \item \textbf{Konumlanma:} EPDK lisanslı CPO \textbf{değil}, \textit{energy management software vendor}.
    \item \textbf{Donanım yatırımı minimum:} 1 lab charger (Vestel Bonus DC ya da kullanılmış EVlink), 1 Power Quality Analyzer, 1 endüstriyel switch, akıllı sayaç. Toplam: 80--150K TL.
\end{itemize}

% ===================================================================
\section{Mevcut Durum — Teknik Değerlendirme}

\subsection{Sahip Olduklarımız}

\begin{tabularx}{\textwidth}{l X}
\toprule
\textbf{Bileşen} & \textbf{Durum} \\
\midrule
Algoritma çekirdeği & 5 controller (UnmanagedController, ManagedController, SRPTController, WaterFillingController, DynamicFairController) \\
Saha senaryoları & 5 (AVM, ofis, otel, hastane, havalimanı) farklı trafo / yük profili \\
Veri modeli & EV, EVModel, ChargingStation, ScenarioConfig, GridConfig, EnvironmentProfile, FleetProfile, ArrivalPattern \\
Veri kaynağı & Sentetik (\texttt{dataset.json}) — \texttt{ArrivalGenerator} + \texttt{BackgroundLoadGenerator} \\
Çıktı & Excel (matris formatlı), matplotlib dashboards (PNG), konsol özet \\
Konfigürasyon & JSON içe/dışa aktarım (\texttt{ScenarioConfig.from\_dict / to\_dict}) \\
\bottomrule
\end{tabularx}

\subsection{Sahip Olmadıklarımız (Kritik Boşluklar)}

\begin{enumerate}
    \item \textbf{Maliyet modeli yok:} Tarife, demand charge, ROI hesabı yok.
    \item \textbf{Gerçek-zamanlı veri yok:} OCPP/Modbus/akıllı sayaç entegrasyonu yok.
    \item \textbf{Yük tahmini yok:} Mevcut \texttt{BackgroundLoadGenerator} sentetik bir model; gerçek tahmin/öğrenme yok.
    \item \textbf{Bulut mimarisi yok:} Mono-script; multi-tenant, API, persistence yok.
    \item \textbf{Müşteri arayüzü yok:} Web dashboard, mobil, alarm yok.
    \item \textbf{Güvenlik yok:} TLS, kimlik doğrulama, KVKK yok.
    \item \textbf{Tahmin/Karar ayrımı yok:} Karar zamanlaması (rolling horizon), ufuk tabanlı optimizasyon yok.
    \item \textbf{Harmonik / güç kalitesi yok:} THD, power factor, dengesizlik analizi yok.
    \item \textbf{Faturalama / ödeme yok:} Şarj seansının parasal karşılığı, fatura, POS yok.
    \item \textbf{PV / ESS entegrasyonu yok:} Solar self-consumption, batarya optimizasyonu yok.
    \item \textbf{Logging / Observability yok:} Üretim seviyesi log yok.
    \item \textbf{CI/CD / Test yok:} Otomatik test, kalite kapısı yok.
\end{enumerate}

\subsection{TRL Seviyesi}

\textbf{Mevcut: TRL 3 (Analitik/Deneysel kavram kanıtlama).} Hedef Faz 1 sonu: TRL 5 (laboratuvar ortamında ilgili teknolojide doğrulama). Hedef Faz 2 sonu: TRL 7 (saha pilotunda gösterim). 12 ay sonu hedef: TRL 8 (gerçek koşullarda tamamlanmış sistem).

% ===================================================================
\section{Boşluk Analizi — 12 Modülün Detaylı Yol Haritası}

Bu bölüm her eksik modülü ayrı ayrı ele alır: \textit{ne lazım, neden lazım, nasıl yaparız, ne kadar sürer, kim yapar.}

% -------------------------------------------------------------------
\subsection{Modül A — Maliyet \& Tarife Modeli}

\subsubsection{Neden Kritik}
Müşteriye ``\%30 tasarruf sağlarım'' diyeceksen tasarrufun para karşılığını hesaplaman lazım. Türkiye'de elektrik tarifesi \textbf{kademeli, zamanlı (TOU), demand-charge'lı}. Müşteri sözleşmesinin ROI cümlesi bu modelden çıkar.

\subsubsection{Türkiye Tarife Yapısı (Sanayi/Ticari, AG-OG)}

\begin{tabularx}{\textwidth}{l l X}
\toprule
\textbf{Bileşen} & \textbf{Birim} & \textbf{Açıklama} \\
\midrule
Aktif enerji & TL/kWh & Üç dilim: Gündüz (06--17), Puant (17--22), Gece (22--06). EPDK her çeyrek günceller. \\
Reaktif enerji & TL/kVArh & Endüktif/kapasitif, oran sınırlarına bağlı (\%20 endüktif, \%15 kapasitif). \\
Güç bedeli & TL/kW.ay & Talep edilen aylık maksimum güç (yalnızca OG müşteri). \\
Dağıtım bedeli & TL/kWh & DSO'ya (BEDAŞ vs.) ödeme. \\
İletim sistemi & TL/kWh & TEİAŞ. \\
Vergi/fonlar & \% & KDV \%20, TRT (\%2 — EV şarjda hariç olabilir, kontrol et), BTV, Enerji Fonu. \\
\bottomrule
\end{tabularx}

\textbf{Önemli:} EV şarjda BTV (\%17) tartışmalı; EPDK ve Hazine duyuruları takip edilmeli.

\subsubsection{Eklenecek Kod Modülü}

\texttt{tariff/} alt-paketi:
\begin{itemize}
    \item \texttt{tariff/tedaş\_2026.json} — güncel tarife verileri.
    \item \texttt{tariff/calculator.py} — \texttt{calculate\_cost(power\_timeseries, tariff, peak\_hours)} → toplam TL.
    \item \texttt{tariff/demand\_charge.py} — aylık 15-dk maksimum güç hesabı.
    \item \texttt{tariff/roi\_simulator.py} — \texttt{simulate\_roi(baseline, optimized, capex)} → geri ödeme süresi.
\end{itemize}

\subsubsection{Süre \& Sahip}
\textbf{Sprint:} 2 hafta. \textbf{Sahip:} CTO/Algoritma. \textbf{Bağımlılık:} EPDK tarife doğrulaması (Operasyon kurucusu).

% -------------------------------------------------------------------
\subsection{Modül B — Gerçek-zamanlı Şebeke Verisi Toplama}

\subsubsection{Soru: ``Anlık veri nereden gelecek?''}

Üç farklı veri kaynağı var, üçü de farklı erişim yolu istiyor:

\begin{tabularx}{\textwidth}{l X X}
\toprule
\textbf{Kaynak} & \textbf{Ne ölçer?} & \textbf{Erişim Yolu} \\
\midrule
Şarj cihazı (EVSE) & Anlık güç, enerji, hata, transaction & \textbf{OCPP 1.6/2.0.1} (websocket/MQTT) \\
Saha akıllı sayaç & Toplam saha gücü, harmonik, PF & \textbf{Modbus TCP/RTU}, \textbf{IEC 61850} (substation), \textbf{DLMS/COSEM} (smart meter standardı) \\
Trafo / OG panel & Trafo yükü, kademe, alarm & Çoğu sahada SCADA yok; akıllı sayaç + akım trafosu (CT) gerekir \\
DSO API'leri & Şebeke kısıtı, talep yanıtı sinyali & \textbf{Türkiye'de henüz açık değil} — BEDAŞ pilot programı dışında \\
EPİAŞ Şeffaflık & Toptan elektrik fiyatı, saatlik üretim/tüketim & Açık API: \texttt{https://seffaflik.epias.com.tr} (OAuth2) \\
\bottomrule
\end{tabularx}

\begin{redbox}
\textbf{Gerçek:} Türkiye'de \textbf{DSO seviyesi açık API yok}. Anlık şebeke yükünü doğrudan BEDAŞ'tan alamazsın. Saha düzeyindeki ölçümü \textbf{kendi akıllı sayacınla} yapacaksın. EPİAŞ sadece toptan piyasa verisi (saatlik granülerlik, ulusal toplam) verir; bireysel saha yükü vermez.
\end{redbox}

\subsubsection{Pratik Strateji}
\begin{enumerate}
    \item \textbf{Saha çıkışına 1 akıllı sayaç} koy — ABB B23, Schneider iEM3000, Janitza UMG (harmonik için).
    \item Sayacın Modbus TCP çıkışını yerel bir IoT gateway'e (Raspberry Pi 4, Siemens IOT2050) bağla.
    \item Gateway $\rightarrow$ MQTT broker (HiveMQ Cloud / EMQX) $\rightarrow$ Bulut ingestion.
    \item Şarj cihazlarının OCPP bağlantısı doğrudan buluta (saha gateway'den geçirme — daha güvenli olur).
\end{enumerate}

\subsubsection{Süre \& Sahip}
\textbf{Sprint:} 4 hafta (lab kurulumu + Modbus okuyucu + MQTT pipeline). \textbf{Sahip:} Embedded/OCPP kurucusu.

% -------------------------------------------------------------------
\subsection{Modül C — OCPP Öğrenme \& Entegrasyon}

\subsubsection{OCPP nedir?}
\textbf{Open Charge Point Protocol} — şarj cihazı (Charge Point) ile merkezi sistem (Central System / CSMS) arasında standartlaştırılmış mesajlaşma protokolü. Open Charge Alliance (OCA) tarafından yönetilir. Bizim için kritik mesajlar:

\begin{tabularx}{\textwidth}{l X}
\toprule
\textbf{Mesaj} & \textbf{Ne işe yarar} \\
\midrule
\texttt{BootNotification} & Charger sisteme katılır; biz onayla, heartbeat aralığı belirleriz. \\
\texttt{Heartbeat} & Charger her N saniyede ``hayatta'' der. \\
\texttt{StartTransaction} / \texttt{StopTransaction} (1.6) / \texttt{TransactionEvent} (2.0) & Şarj seansı başla/bitti. \\
\texttt{MeterValues} & Anlık güç, enerji, gerilim, akım. \\
\texttt{StatusNotification} & Soket durumu (Available, Charging, Faulted...). \\
\texttt{SetChargingProfile} & \textbf{Bizim güç limiti gönderdiğimiz mesaj — hayati.} \\
\texttt{RemoteStartTransaction} / \texttt{RemoteStopTransaction} & Uzaktan şarj başlat/durdur. \\
\bottomrule
\end{tabularx}

\subsubsection{Öğrenme Yolu (4 Hafta)}

\textbf{Hafta 1 — Standart oku:}
\begin{itemize}
    \item OCPP 1.6 Edition 2 spec PDF (yaklaşık 90 sayfa).
    \item Open Charge Alliance üyeliği (\$1500/yıl) — lab kullanım için zorunlu değil ama sertifikasyon için gerekli.
    \item ``OCPP Smart Charging'' beyaz kitap.
\end{itemize}

\textbf{Hafta 2 — Lab simülasyon:}
\begin{itemize}
    \item \texttt{mobilityhouse/ocpp} (Python kütüphanesi) — açık kaynak, en yaygın.
    \item \texttt{SteVe} (Java backend) — hazır CSMS, kod okumak için iyi.
    \item Sahte charger simülatörü: \texttt{ocpp-cp-py} veya MAEVe.
    \item Lokalde 1 ``CSMS'' + 5 ``charger'' simülatörü çalıştır, mesaj akışını gör.
\end{itemize}

\textbf{Hafta 3 — Gerçek charger lab testi:}
\begin{itemize}
    \item Vestel Bonus DC veya 2. el ABB Terra alın (kira da olur).
    \item Charger'ı internete (kendi VPN ya da DDNS ile) bağla.
    \item OCPP konfigürasyonunu CSMS URL'ine yönlendir.
    \item İlk \texttt{BootNotification} alındığı an sevin.
\end{itemize}

\textbf{Hafta 4 — Smart Charging:}
\begin{itemize}
    \item \texttt{SetChargingProfile} ile statik 11 kW limit gönder.
    \item Sonra dinamik limit: simülasyondaki \texttt{DynamicFairController}'dan çıkan limit\_kw'yi gerçek charger'a uygula.
    \item Validasyon: charger gerçekten limit'e uyuyor mu? \texttt{MeterValues}'dan kontrol et.
\end{itemize}

\subsubsection{Charger Üreticileri Türkiye Listesi}

\begin{tabularx}{\textwidth}{l l X}
\toprule
\textbf{Üretici} & \textbf{Model Örneği} & \textbf{OCPP Notu} \\
\midrule
Vestel & Bonus AC22, Bonus DC60--180 & 1.6J, Smart Charging mevcut \\
Phihong & DC fast & 1.6J, 2.0 yol haritası \\
ABB & Terra AC, Terra DC & 1.6J + 2.0.1 (yeni firmware) \\
EnerjiSA & Eşarj cihazları & 1.6J \\
ZES & Kendi cihazları + 3rd-party & 1.6J \\
Schneider & EVlink Pro AC, EVlink DC & 1.6J + 2.0.1 \\
Wallbox (İspanya) & Pulsar, Quasar & 1.6J + 2.0 \\
Charpoint (US) & CT4000 & 1.6J — daha kapalı ekosistem \\
\bottomrule
\end{tabularx}

\textbf{Tavsiye:} İlk pilotta Vestel veya ABB Terra. Sebep: Türkiye servis ağı + iyi OCPP destek + makul fiyat.

\subsubsection{Süre \& Sahip}
\textbf{Sprint:} 4 hafta öğrenme + 4 hafta production-ready entegrasyon. \textbf{Sahip:} Embedded/OCPP kurucusu birincil; CTO destek.

% -------------------------------------------------------------------
\subsection{Modül D — Lokal Yük Tahmini (ML/AI)}

\subsubsection{Soru: ``Model mi eğiteceğiz, veri seti var mı?''}

\textbf{Cevap: Evet, eğiteceğiz. Kendi veriniz olana kadar transfer learning + public veri.}

\subsubsection{Tahmin Edilecek Şey Ne?}

İki ayrı tahmin problemi var:
\begin{enumerate}
    \item \textbf{Saha baz yükü tahmini} (background load forecast) --- AVM/otel/fabrikanın sonraki 1-24 saat tüketimi. Bu mevcut sim'deki \texttt{BackgroundLoadGenerator}'ın yerini alır.
    \item \textbf{EV gelişi/SoC tahmini} (EV arrival forecast) --- sonraki saatlerde kaç araç gelecek, tahmini SoC dağılımı. Bu \texttt{ArrivalGenerator}'ın yerini alır.
\end{enumerate}

\subsubsection{Veri Seti Stratejisi}

\textbf{Public veri setleri (eğitim için ön-prior):}

\begin{tabularx}{\textwidth}{l X l}
\toprule
\textbf{Veri seti} & \textbf{İçerik} & \textbf{Lisans} \\
\midrule
ENTSO-E Transparency & TR dahil 30+ ülke saatlik tüketim/üretim & Açık \\
EPİAŞ Şeffaflık & TR saatlik tüketim, üretim, fiyat & Açık (kayıt) \\
UCI ML Repo --- ``Electricity Load Diagrams'' & Portekiz 370 müşteri 15-dk yük & Açık \\
UMass Smart$\ast$ Dataset & ABD bina + apartman & Akademik \\
Pecan Street Dataport & ABD residence (EV charging dahil) & Akademik (\$\$ ticari) \\
ACN-Data (Caltech) & ABD üniversite kampüsü EV şarj & Akademik \\
REFIT (UK) & 20 ev, alet seviyesi, EV yok & Akademik \\
GREEND (AT/IT) & 8 ev & Akademik \\
ChargeCar (CMU) & EV sürüş profilleri (şarj yok) & Akademik \\
OpenChargeMap & Şarj noktası lokasyonları (yük yok) & Açık \\
NREL EV Charging Data & ABD CPO verileri (anonim) & Açık \\
\bottomrule
\end{tabularx}

\textbf{Pratik:} Public veri Türkiye gerçeği için \textbf{yetersiz}. Strateji:
\begin{enumerate}
    \item \textbf{Pre-training:} Public veride güçlü bir \textit{baseline forecaster} yetiştir (özellikle EPİAŞ saatlik TR ulusal tüketim).
    \item \textbf{Pilot toplama:} İlk friendly pilotta 3 ay boyunca dakikalık veri biriktir (akıllı sayaç).
    \item \textbf{Fine-tuning:} Her saha için son N hafta veriyle modeli site-specific yap (transfer learning).
    \item \textbf{Online learning:} Üretimde günlük kayan pencerede yeniden eğit (retrain weekly).
\end{enumerate}

\subsubsection{Model Seçimi (Aşamalı)}

\begin{tabularx}{\textwidth}{l l l X}
\toprule
\textbf{Aşama} & \textbf{Model} & \textbf{Tahmin Ufku} & \textbf{Not} \\
\midrule
V1 (Hafta 1-2) & Naive: ``geçen haftanın aynı saati'' & 1-24 sa & Baseline; akıllıca tasarlanırsa şaşırtıcı iyi. \\
V1 (Hafta 3-4) & SARIMA / Prophet & 1-24 sa & Açıklanabilir, kurumsal müşteri sever. \\
V2 (Ay 2-4) & XGBoost / LightGBM (lag features + takvim + hava) & 1-24 sa & Ana iş yükü modeli; üretim baseline. \\
V3 (Ay 4-7) & LSTM / GRU & 1-72 sa & Daha uzun ufuk; karmaşıklık artar. \\
V4 (Ay 7+) & Temporal Fusion Transformer (TFT) / NeuralProphet / Chronos (Amazon) & 1-168 sa & State-of-the-art; yatırımcı toplantısında parlatıcı. \\
\bottomrule
\end{tabularx}

\textbf{Tavsiye:} V2 (XGBoost) üretime ilk koy. V3/V4 araştırma + diferansiyatör. Kurumsal müşteri ``Transformer'' kelimesini sever ama gerçekte XGBoost yeterince iyidir.

\subsubsection{Feature Engineering — Önemli Özellikler}

\begin{itemize}
    \item \textbf{Zaman:} hour-of-day, day-of-week, day-of-month, month, is\_weekend, is\_holiday (TR resmi tatil).
    \item \textbf{Lag:} t-1h, t-24h, t-168h (geçen hafta aynı saat). Mavi-altın bilgi.
    \item \textbf{Rolling stats:} son 1h/24h/7d ortalama, std.
    \item \textbf{Hava:} sıcaklık, nem, ışık (MGM açık API), tahmin (OpenWeather).
    \item \textbf{Takvim:} okul tatili, ramazan/bayram (AVM yükü değişir), maç günü (otopark dolar).
    \item \textbf{Yerel olay:} mağaza açılış/kapanış, bina işletmesi vardiyaları.
    \item \textbf{EV özel:} son 24 saatte gelen araç sayısı (otokorelasyon güçlü).
\end{itemize}

\subsubsection{Validasyon \& Metrik}
\begin{itemize}
    \item Walk-forward CV (zaman serileri için zorunlu, KFold yanıltıcı).
    \item Metrik: MAE (kW), MAPE (\%), RMSE.
    \item Pinball loss --- quantile forecasting (P10/P50/P90) için. Karar verme uncertainty-aware olur.
    \item Kabul kriteri pilotta: 1-saat-ileri MAPE $<$\%10, 24-saat-ileri MAPE $<$\%18.
\end{itemize}

\subsubsection{Süre \& Sahip}
\textbf{Sprint:} V1 baseline (2 hafta), V2 XGBoost (3 hafta), V3 LSTM (4 hafta + araştırma payı). \textbf{Sahip:} CTO/Algoritma birincil; akademik partner (İTÜ EE / ODTÜ EE) tez/intern desteği.

% -------------------------------------------------------------------
\subsection{Modül E — Optimizasyon Algoritması (Karar)}

\subsubsection{Mevcut Durum}
\texttt{simulation\_v3.py} 5 algoritma:
\begin{itemize}
    \item \textbf{UnmanagedController} --- baseline (yalın).
    \item \textbf{ManagedController} --- starvation + bitiş yakınlığı priorite.
    \item \textbf{SRPTController} --- en az kalan iş süresi.
    \item \textbf{WaterFillingController} --- adil paylaşım.
    \item \textbf{DynamicFairController} --- ağırlıklı urgency.
\end{itemize}
Hepsi \textbf{greedy / heuristik}. Anlık karar, ileriyi görmüyor.

\subsubsection{Eklenecek (Faz 2-3)}

\textbf{1. Rolling Horizon MILP/QP optimizasyonu:}
\begin{itemize}
    \item Karar değişkeni: $P_{i,t}$ — araç $i$, dakika $t$ gücü.
    \item Kısıt:
        \begin{itemize}
            \item $P_{i,t} \le P_{i,\max}$ (charger × araç limiti).
            \item $\sum_i P_{i,t} + B_t \le L_t$ (toplam ≤ trafo limiti, $B$ baz yük tahmini, $L$ trafo).
            \item $\sum_t P_{i,t} \cdot \Delta t \ge E_i$ (her araç enerjisini al).
        \end{itemize}
    \item Hedef: $\min \sum_t c_t \cdot \big(\sum_i P_{i,t} + B_t\big) + \lambda \sum_i \text{wait}_i$ (maliyet + bekleme cezası).
    \item Çözücü: CVXPY (Python, açık), OR-Tools (Google), Gurobi (lisans gerekli, hızlı).
    \item Ufuk: 4--12 saat ileri, her 5--15 dakikada yeniden çöz.
\end{itemize}

\textbf{2. Stokastik / Robust Optimizasyon:}
\begin{itemize}
    \item Tahmin belirsizliği var (P10/P90).
    \item Scenario-based stochastic LP: birkaç ``senaryo'' örnekle, her birini optimize et, beklenen değeri minimize et.
    \item Daha güvenli ama daha pahalı; faz 3.
\end{itemize}

\textbf{3. Reinforcement Learning (Faz 4 araştırma):}
\begin{itemize}
    \item Sim ortamı zaten var (\texttt{simulation\_v3.py}). PPO/SAC ajanı eğit.
    \item Avantaj: implicit objective öğrenimi.
    \item Dezavantaj: kara kutu — kurumsal müşteri sevmez.
    \item Akademik ortaklıkla (TÜBİTAK 1505 ARDEB) tez konusu.
\end{itemize}

\subsubsection{Süre \& Sahip}
MILP rolling horizon: 6 hafta. Stokastik: ek 4 hafta. RL: araştırma — 6+ ay. \textbf{Sahip:} CTO/Algoritma.

% -------------------------------------------------------------------
\subsection{Modül F — Güç Kalitesi \& Harmonik (Faz 4 Diferansiyatör)}

\subsubsection{Neden Önemli}

DC hızlı şarj cihazları \textbf{ciddi harmonik üretir} (özellikle eski 6-pulse rectifier'lı modeller). Saha trafosu / sayacı bunu görür ve:
\begin{itemize}
    \item DSO (BEDAŞ) reaktif/harmonik cezası kesebilir (IEEE 519 TR uyarlaması).
    \item Trafo gerçek yükü (THD dahil) görünür yükten yüksektir → anlık aşım.
    \item Diğer cihazlara zarar verebilir (motor, kompresör).
\end{itemize}

\textbf{Ticari fırsat:} ``Sadece güç limiti değil, harmonik-aware load balancing yapıyoruz'' demek diferansiyatör.

\subsubsection{Standartlar}

\begin{itemize}
    \item \textbf{IEEE 519-2014} --- harmonic limits (THD\_V, THD\_I limitleri).
    \item \textbf{IEC 61000-3-12} --- 16--75 A cihaz harmonik.
    \item \textbf{IEC 61000-4-7} --- ölçüm yöntemi.
    \item \textbf{TS EN 50160} --- şebeke gerilim kalitesi (TR).
\end{itemize}

\subsubsection{Donanım}

\begin{itemize}
    \item \textbf{Power Quality Analyzer} --- Janitza UMG 605 (Modbus + harmonic up to 50.), Fluke 437-II (lab).
    \item \textbf{Akım trafoları (CT)} --- Saha çıkışı + her charger besleme hattı için.
    \item \textbf{Veri yolu:} Modbus TCP $\rightarrow$ kendi gateway $\rightarrow$ MQTT.
\end{itemize}

\subsubsection{Yazılım}

\begin{itemize}
    \item Anlık FFT analizi --- gerek var mı? Hayır, Janitza zaten FFT yapar; biz harmonik ölçümleri okuruz.
    \item Algoritma: ``charger önceliklendirmede THD'yi de objective'e ekle''. Mevcut algoritmaya ek skor.
    \item Alarm: THD\_V $>$ \%5 → otomatik kısma.
\end{itemize}

\subsubsection{Süre \& Sahip}
\textbf{Faz 4 araştırma — Ay 9+.} 6 hafta. \textbf{Sahip:} Embedded + akademik partner (İTÜ Güç Sistemleri, Yıldız EE).

\begin{notebox}
\textbf{Stratejik not:} Harmonik analizini Faz 1--2'de yapma, ama \textbf{satış pitch'inde bahsediyoruz} ki ``ürün yol haritamızda var, bizden bunu istersen yapacağız.'' Yatırımcı / büyük kurumsal müşteri için diferansiyatör.
\end{notebox}

% -------------------------------------------------------------------
\subsection{Modül G — Bulut Mimarisi (Multi-tenant SaaS)}

\subsubsection{Mevcut → Hedef}

\textbf{Mevcut:} \texttt{python simulation\_v3.py} → 1 makine, 1 senaryo.

\textbf{Hedef:} N müşteri × M saha × K charger × T zaman = bulut tabanlı multi-tenant SaaS.

\subsubsection{Mimari Katmanlar}

\begin{lstlisting}[basicstyle=\small\ttfamily, frame=single]
┌─────────────────────────────────────────────────────┐
│ Web Dashboard (Next.js, Vercel)                     │
│ Mobile (React Native, opsiyonel — Faz 3)            │
├─────────────────────────────────────────────────────┤
│ API Gateway (Kong/Traefik) + WAF                    │
│   - REST/GraphQL                                    │
│   - Rate limit, auth (Auth0 / Keycloak)             │
├─────────────────────────────────────────────────────┤
│ Application Servers (FastAPI)                       │
│   - Tenant/Site/Charger CRUD                        │
│   - Algoritma çekirdeği (simulation_v3 → port)     │
│   - Tariff/ROI servisi                              │
├─────────────────────────────────────────────────────┤
│ OCPP Gateway (Python, mobilityhouse/ocpp)           │
│   - WebSocket connection per charger               │
│   - Telemetry → Kafka                              │
│   - Command (SetChargingProfile) ← API              │
├─────────────────────────────────────────────────────┤
│ Real-time Optimization Worker                       │
│   - Her saha için ayrı container                    │
│   - 30 sn / 1 dk cycle: forecast → optimize → push │
├─────────────────────────────────────────────────────┤
│ ML Inference Service (load forecast)                │
│   - FastAPI + ONNX runtime / TF Serving             │
├─────────────────────────────────────────────────────┤
│ Streaming (Kafka / NATS) — telemetry pipeline       │
│ TimescaleDB (zaman serisi)                          │
│ PostgreSQL (durum, kullanıcı, tenant)               │
│ Redis (cache, distributed lock)                     │
│ S3-compatible object store (rapor, ML artifacts)    │
├─────────────────────────────────────────────────────┤
│ Observability: Grafana + Prometheus + Loki          │
│ Errors: Sentry                                      │
│ Logs: ELK / Loki                                    │
└─────────────────────────────────────────────────────┘
\end{lstlisting}

\subsubsection{Cloud Provider Kararı}

\begin{tabularx}{\textwidth}{l X X}
\toprule
\textbf{Sağlayıcı} & \textbf{Avantaj} & \textbf{Dezavantaj} \\
\midrule
AWS (Frankfurt) & Olgun, geniş hizmet, KVKK uyumlu (DPA imzalanır) & Maliyet (TL düşüşünde acı) \\
Google Cloud (Frankfurt) & İyi ML/AI servisleri & TR pazarda az tanınır \\
Azure (Frankfurt) & Kurumsal müşteri sever & Maliyet \\
Turkcell BulutTürk & Yerli, KVKK kolay, kamu satışında avantaj & Hizmet derinliği daha az \\
Türk Telekom Bulut & Yerli + KOBİ desteği & ML kit yok \\
\bottomrule
\end{tabularx}

\textbf{Tavsiye:} Faz 1--2'de \textbf{Hetzner Cloud (Almanya, ucuz)} + \textbf{Cloudflare (CDN/WAF)}. Maliyet kontrolü için. Faz 3'te kurumsal/devlet için \textbf{Turkcell BulutTürk}'e migration opsiyonu.

\subsubsection{Multi-tenancy Yaklaşımı}
\begin{itemize}
    \item \textbf{Schema-per-tenant} (Postgres). Çoğu kurumsal müşteri ``benim verim ayrı'' diye ısrarcı; bu mimari karar onu çözer.
    \item Row-level security (RLS) ile tek-DB yaklaşımı erken evrede daha hızlı; faz 2'de schema-per-tenant'a geç.
\end{itemize}

\subsubsection{Süre \& Sahip}
\textbf{V1 MVP backend:} 4 hafta (CTO + Embedded). \textbf{V2 prod-ready:} +8 hafta. \textbf{Sahip:} CTO birincil.

% -------------------------------------------------------------------
\subsection{Modül H — Güvenlik \& Uyum}

\subsubsection{Boyutlar}

\begin{tabularx}{\textwidth}{l X}
\toprule
\textbf{Boyut} & \textbf{Aksiyon} \\
\midrule
TLS / Şifreleme & Tüm OCPP bağlantısı WSS (TLS 1.2+); cert-manager + Let's Encrypt. \\
Auth & OAuth2 + JWT; refresh token; MFA (TOTP) admin için. \\
KVKK & Aydınlatma metni, VERBİS sicili, ROPA, veri saklama politikası, silme talebi akışı. \\
ISO 27001 & Faz 2 sonu hedef. Asistan: 6 ay süre, dış denetim. \\
IEC 62443 & Endüstriyel kontrol güvenliği — kurumsal müşteri sorabilir. Faz 3. \\
Penetration test & Yılda 1 kez (3rd party). PreSeed sonrası gerçekçi. \\
Secrets mgmt & HashiCorp Vault veya cloud KMS; environment varlığında düz parola \textbf{yok}. \\
RBAC & Kullanıcı: admin, operatör, salt-okuyan. Tenant izolasyonu. \\
Audit log & Tüm yazma işlemi loglanır, 1 yıl saklanır. \\
\bottomrule
\end{tabularx}

\subsubsection{KVKK Detayı (atlanırsa cezası ağır)}

\begin{itemize}
    \item Şarj seansı = kişisel veri (kullanıcı RFID/araç plaka eşleşmesi).
    \item \textbf{Veri Sorumlusu sicili (VERBİS)} kayıt zorunlu (yıllık ciro 25M TL üzerinde olsun olmasın, 50+ çalışan üzerinde olmasın gerek yok ama 100 müşteri verisi varsa kaydet).
    \item \textbf{Aydınlatma metni} her müşteri sözleşmesinde ek olmalı.
    \item \textbf{ROPA (Veri İşleme Faaliyetleri Kaydı)} — iç dosya, denetimde sorulur.
    \item \textbf{Veri ihlali bildirimi} 72 saat içinde KVKK Kurulu'na.
    \item \textbf{Yurt dışı aktarım} (AWS Frankfurt) — DPA imzala, açık rıza al veya KVKK'nın belirlediği uygun ülke listesine uy.
\end{itemize}

\subsubsection{Süre \& Sahip}
KVKK uyum: 2 hafta + yıllık güncelleme. ISO 27001: 6 ay (faz 2). Pentest: yılda 1 (50--150K TL dış maliyet). \textbf{Sahip:} CEO/Operasyon (hukuk dış kaynak).

% -------------------------------------------------------------------
\subsection{Modül I — PV / ESS Entegrasyonu (Faz 3 Diferansiyatör)}

\subsubsection{Niye Önemli}
Saha sahibi çatıya PV koyduysa veya batarya BESS aldıysa, ``EV şarjını PV'den karşıla, fazlası varsa şebekeye sat, eksiği batarya'dan al'' bir gold mine. Türkiye'de:
\begin{itemize}
    \item Çatı GES (lisanssız) yaygınlaştı, AVM/fabrika çatıları doluyor.
    \item BESS (Tesla Megapack, BYD, Sungrow) henüz pahalı ama 2026--2027 cazip olacak.
    \item EPDK BESS yönetmeliği (2023) yayımlandı; \textbf{aggregator} olabilir.
\end{itemize}

\subsubsection{Algoritma Genişlemesi}

Bizim yük dengeleme algoritmamızın objective'ine 3 yeni terim ekle:
\begin{enumerate}
    \item \textbf{PV self-consumption} maksimize et (PV anlık $\rightarrow$ EV).
    \item \textbf{BESS scheduling} --- gece ucuz tarifede şarj et, puantta deşarj.
    \item \textbf{Şebeke export} --- YEKDEM dışı satış (sınırlı, AB modelinde değil).
\end{enumerate}

\subsubsection{Süre \& Sahip}
Faz 3, Ay 8+. 6--8 hafta. \textbf{Sahip:} CTO + akademik partner.

% -------------------------------------------------------------------
\subsection{Modül J — V2G \& V2H (Faz 4)}

V2G (Vehicle-to-Grid) — aracın bataryasından şebekeye / binaya geri elektrik akışı. Türkiye'de:
\begin{itemize}
    \item Henüz yaygın değil.
    \item TOGG T10X V2L (vehicle-to-load) yapabilir, V2G şu an yok.
    \item ISO 15118-20 + CCS2 ile teorik mümkün.
    \item Bizim için 12--18 ayda \textbf{pitch maddesi}, ürün maddesi değil.
\end{itemize}
\textbf{Tavsiye:} Mimaride ``hazır'' yap (bidirectional power flow desteği), gerçek implementasyon Faz 4 sonrası.

% -------------------------------------------------------------------
\subsection{Modül K — Müşteri Arayüzleri}

\subsubsection{Operatör Dashboard (web)}

Kim kullanır: AVM teknik müdürü, lojistik filo amiri, otel facility manager.

Ana ekranlar:
\begin{itemize}
    \item \textbf{Anlık durum:} saha güç (kW), trafo limiti, charger durumu (gantt-vari).
    \item \textbf{Tahmin:} sonraki 24 saat baz yük + tahmini EV gelişi.
    \item \textbf{Olay/Alarm:} aşım, charger arıza, OCPP kopması.
    \item \textbf{Rapor:} aylık tasarruf, fatura öncesi/sonrası, charger başına servis.
    \item \textbf{Konfigürasyon:} algoritma seçimi, priorite kuralları, alarm eşikleri.
\end{itemize}

\subsubsection{Saha Mobile (Faz 3 sonra)}
Sürücü için: rezervasyon, durum, ödeme. \textbf{Şu an gerek yok} — Voltrun/ZES uygulamaları zaten var, biz onların altında çalışırız.

\subsubsection{API}
Üçüncü parti integrasyon (CPO platformları, BMS sistemleri, Energy Management Systems). REST + Webhook.

\subsubsection{Süre \& Sahip}
Operatör dashboard MVP: 4 hafta (CEO/Operasyon + frontend hire/freelancer). Mobile: Faz 3.

% -------------------------------------------------------------------
\subsection{Modül L — Faturalama / Ödeme}

\subsubsection{Soru: Müşteriden nasıl para alacağız?}

\textbf{Bizim model SaaS:}
\begin{itemize}
    \item Setup fee (30--80K TL): banka havalesi, e-fatura.
    \item Aylık SaaS (charger × ay × birim fiyat): otomatik fatura, vade 30 gün.
    \item Sözleşme: KKTC/TR muhasebe, e-fatura/e-arşiv (GİB).
\end{itemize}

\textbf{Saha sahibi araç sürücüsünden nasıl alacak?} Bizim sorunumuz değil ama mimari etki var:
\begin{itemize}
    \item Eğer saha bir RFID/QR sistemi kullanıyorsa, biz \texttt{user\_id}'yi seans kaydında tutarız.
    \item Saha CPO'ya bağlıysa (Voltrun gibi), CPO'nun billing'i kendi yapar; biz onun altyapısıyız.
\end{itemize}

\subsubsection{Faturalama Yazılımı (Bizim için)}
\begin{itemize}
    \item Logo / Mikro / Paraşüt — küçük şirket için yeterli.
    \item Stripe Türkiye ya da iyzico — kart ödeme (gelecek).
    \item TR e-fatura entegrasyonu zorunlu (müşteri ciro $>$ 5M TL).
\end{itemize}

% -------------------------------------------------------------------
\subsection{Modül M — Atlanmış / Daha Az Konuşulmuş 25 Konu}

\begin{enumerate}
    \item \textbf{Batarya degradasyonu} --- yüksek hızlı (HPC) şarj batarya ömrünü azaltır. Müşteriye söylenecek bir trade-off; ``adil'' algoritma slow-charge'ı tercih edebilir. Akademik makale potansiyeli.
    \item \textbf{SLA tanımı} --- kurumsal müşteri ``\%99.5 uptime garantisi'' ister. Cezası kontratta. Mimari etkiler bunun üzerine kurulur.
    \item \textbf{Disaster Recovery (DR) / Backup site} --- Türkiye'de deprem riski; AB region (Frankfurt) hot standby. Faz 2 sonu.
    \item \textbf{Charger Firmware Update Yönetimi} --- OCPP üzerinden firmware push. \texttt{UpdateFirmware} mesajı. Saha-saha firmware versiyon takibi şart.
    \item \textbf{Data Governance / Veri Sahipliği} --- ``müşteri veri sahibi, biz veri işleyen''. Sözleşmede netleştir. Anonimleştirilmiş istatistik için ek lisans.
    \item \textbf{Anomaly Detection} --- charger sahteciliği, RFID kötüye kullanım, fraud. Faz 2.
    \item \textbf{Logging \& Observability} --- prod'da Grafana panel, alarm. SLO/SLI tanımları.
    \item \textbf{Multi-language UI} --- TR ana, EN second. MENA için Arapça orta vade. i18n stack erkenden kur.
    \item \textbf{Edge computing} --- bazı sahalar (ADM, hastane) ``veri dışarı çıkmasın'' der. On-prem bir node opsiyonu lazım. Faz 3.
    \item \textbf{Şarj Cihazı Offline Davranışı} --- internet kopunca charger ne yapar? OCPP 1.6 \texttt{LocalAuthList} ve \texttt{ChargingProfilesPersistent} ile graceful degradation. Müşteriye anlatılması gereken.
    \item \textbf{Sözleşme Şablonları} --- pilot, ticari, OEM ortak. Hukuki danışman ile (Esin Avukatlık veya Pekin Bayar Türkiye) standartlaştır.
    \item \textbf{Patent Stratejisi} --- TR/EPO algoritma patenti zor; ama \textit{trade secret} + müşteri sözleşmesinde IP koruma. Belki ``yöntem patenti'' ABD'de mümkün.
    \item \textbf{Marka / Domain / Trademark} --- isim seç, .com + .com.tr al, TPMK marka başvurusu (3K TL, 6 ay).
    \item \textbf{Sektör Konferansları} --- ICCI (Mart, İstanbul), Solarex (Nisan, İstanbul), EVER Monaco, ChargeIN Europe. Demo standı.
    \item \textbf{Sektör Dergileri / PR} --- Enerji Dergisi, Elektrik Mühendisliği. Düzenli yazı + case study yayını = inbound lead.
    \item \textbf{Akademik Yayın} --- IEEE Trans. Smart Grid, Applied Energy. Diferansiyel olur, yatırımcı sever.
    \item \textbf{Open-source Stratejisi} --- OCPP gateway'i open source yap (Apache 2.0). Algoritma core'u kapalı. Developer mindshare.
    \item \textbf{Müşteri Başarı (Customer Success)} --- pilot başına dedicated CSM (Customer Success Manager). Faz 2 hire.
    \item \textbf{Sigorta} --- kurumsal müşteri ``Cyber Liability Insurance + Professional Indemnity'' ister. AXA, Allianz, Türkiye'de Anadolu Sigorta.
    \item \textbf{ISO 9001 (Kalite)} --- devlet ihalesinde sorulur. Faz 3 nice-to-have.
    \item \textbf{ISO 14001 (Çevre)} --- yine devlet ihalesinde. Faz 3.
    \item \textbf{TS 13298 / e-İmza / KEP} --- kurumsal müşteri sözleşmelerinde KEP. Erkenden al (Türk Telekom, E-Güven).
    \item \textbf{Talep Yanıtı (Demand Response)} --- Türkiye'de henüz EPİAŞ üzerinden bireysel saha katılımı sınırlı. Aggregator rolü Faz 3.
    \item \textbf{Yedek Şirket Yapısı} --- TR Ltd Şti + Estonya OÜ (e-Residency) iki katmanlı yapı; fonlama ve uluslararası satış için kolaylık. Hukuki danışman gerekli.
    \item \textbf{Founder vesting + ESOP havuzu} --- ESOP \%10--15 erken evrede, ileride hire için kritik. Founders' Agreement gün 1.
\end{enumerate}

% ===================================================================
\section{Veri Stratejisi (Detay)}

\subsection{İhtiyacımız Olan Veriler}

\begin{tabularx}{\textwidth}{l X l}
\toprule
\textbf{Veri} & \textbf{Niye} & \textbf{Kaynak} \\
\midrule
Saha toplam yük (1-dk) & Tahmin modeli, anlık karar & Akıllı sayaç (saha çıkışı) \\
Charger telemetri (1-sn / 1-dk) & Anlık güç, transaction & OCPP MeterValues \\
EV bilgisi (model, batarya, SoC) & SoC kestirimi & OCPP TransactionEvent + 15118 \\
Tarife verisi (saatlik) & Maliyet hesabı, optimizasyon & EPDK + EPİAŞ \\
Hava verisi (saatlik) & Yük tahmini özniteliği & MGM açık API, OpenWeather \\
Takvim (TR tatil) & Yük tahmini & Statik tablo + ramazan/bayram lib \\
Ulusal saatlik tüketim & Pre-training prior & EPİAŞ Şeffaflık \\
Yerel olay (maç, tatil) & İleri yük tahmini & Manuel etiketleme + scrape \\
\bottomrule
\end{tabularx}

\subsection{Pilot Veri Toplama Planı}

\begin{enumerate}
    \item Pilot başında saha çıkışına Janitza UMG kur (4 hafta lab + saha kurulum).
    \item İlk 4 hafta sadece ``shadow mode'' --- kontrol etme, sadece veri biriktir.
    \item 4--8 hafta: yük tahmini eğit (200K+ örnek).
    \item 8--16 hafta: kontrol modu aç, A/B test (yarısı müdahalesiz, yarısı algoritma).
\end{enumerate}

\subsection{Veri Yönetişimi}

\begin{itemize}
    \item Müşteri veri sahibi; biz veri işleyen.
    \item Anonimleştirilmiş aggregate veri için açık rıza (sözleşmede checkbox).
    \item Veri saklama: 24 ay default, müşteri talebi ile silme.
    \item Bulut: AWS Frankfurt + DPA imzası.
\end{itemize}

% ===================================================================
\section{Mimari Evrim --- Sürüm Yol Haritası}

\begin{tabularx}{\textwidth}{l l X}
\toprule
\textbf{Sürüm} & \textbf{Hedef} & \textbf{Ana içerik} \\
\midrule
V0 (mevcut) & Lab POC & \texttt{simulation\_v3.py} mono-script. \\
V1 (Ay 2--4) & MVP & FastAPI backend + 1 charger OCPP entegrasyonu + read-only dashboard + sentetik veri. \\
V2 (Ay 4--6) & Pilot-ready & Multi-tenant, gerçek tarife, ML baseline (XGBoost), tahmin + karar entegre, friendly pilot canlı. \\
V3 (Ay 6--9) & Ücretli pilot & 3--5 saha, Grafana izleme, KVKK/güvenlik full, alarm/uyarı, müşteri raporu. \\
V4 (Ay 9--12) & Scale-ready & DR, multi-region, OCPP 2.0.1, demand response sinyali, PV/ESS proto. \\
V5 (Ay 12+) & Diferansiyatör & Harmonik analiz, V2G hazır, RL araştırma. \\
\bottomrule
\end{tabularx}

% ===================================================================
\section{12 Aylık Birleşik Sprint Planı (Hafta Hafta)}

\textbf{Notasyon:} {\color{blue!50!black}T} = Teknik, {\color{red!50!black}İ} = İdari/İş, {\color{green!50!black}P} = Pazarlama/Satış.

\subsection*{Ay 1 — Validasyon \& Hazırlık}

\textbf{Hafta 1:}
\begin{itemize}
    \item {\color{red!50!black}İ} Founders' Agreement, vesting, ESOP havuzu, şirket türü kararı.
    \item {\color{red!50!black}İ} Domain + Linkedin + e-posta + GitHub org.
    \item {\color{red!50!black}İ} Hukuki danışman ilk toplantı (KVKK + EPDK + sözleşme şablonları).
    \item {\color{green!50!black}P} 30 müşteri görüşme listesi, outreach mesajları.
\end{itemize}

\textbf{Hafta 2--3:}
\begin{itemize}
    \item {\color{green!50!black}P} 30 customer discovery (3 kurucu × 10).
    \item {\color{blue!50!black}T} OCPP 1.6 spec oku; mobilityhouse/ocpp lokal kur.
    \item {\color{red!50!black}İ} TÜBİTAK 1512 BiGG başvuru paketi başla.
\end{itemize}

\textbf{Hafta 4:}
\begin{itemize}
    \item {\color{green!50!black}P} Pain matrisi, 3 LOI hedef.
    \item {\color{blue!50!black}T} Lab simülasyonda 5 sahte charger ile CSMS bağlantısı.
    \item {\color{red!50!black}İ} BiGG son kontrol + başvur.
\end{itemize}

\subsection*{Ay 2 — Donanım Lab + V1 Backend}

\textbf{Hafta 5:}
\begin{itemize}
    \item {\color{blue!50!black}T} Vestel Bonus DC sipariş (3--6 hafta termin), Janitza UMG sipariş.
    \item {\color{blue!50!black}T} FastAPI proje skeleton + Postgres + Redis. Tenant/Site/Charger CRUD.
    \item {\color{blue!50!black}T} \texttt{simulation\_v3.py} algoritmalarını \texttt{evlb/algorithms/} paketi olarak refactor (production reuse).
\end{itemize}

\textbf{Hafta 6--7:}
\begin{itemize}
    \item {\color{blue!50!black}T} OCPP gateway (mobilityhouse/ocpp) entegre. WSS endpoint.
    \item {\color{blue!50!black}T} Telemetry → TimescaleDB.
    \item {\color{blue!50!black}T} Tarife modülü (Modül A) — Türkiye TEDAŞ tarife jsonu.
    \item {\color{green!50!black}P} 2 friendly pilot LOI imza.
\end{itemize}

\textbf{Hafta 8:}
\begin{itemize}
    \item {\color{blue!50!black}T} Lab kurulumu: Vestel Bonus + Janitza + IoT gateway.
    \item {\color{blue!50!black}T} İlk gerçek charger \texttt{BootNotification} alındı --- şampanya patlat.
\end{itemize}

\subsection*{Ay 3 — Smart Charging + ROI Modülü}

\textbf{Hafta 9--10:}
\begin{itemize}
    \item {\color{blue!50!black}T} \texttt{SetChargingProfile} ile statik limit gönder. Doğrulama.
    \item {\color{blue!50!black}T} ROI calculator — pilot teklifi için canlı.
    \item {\color{green!50!black}P} İlk friendly pilot saha kurulum (sözleşme imza).
\end{itemize}

\textbf{Hafta 11--12:}
\begin{itemize}
    \item {\color{blue!50!black}T} DynamicFairController production-ready (1-dk cycle).
    \item {\color{blue!50!black}T} Dashboard read-only — anlık güç + transaction listesi.
    \item {\color{red!50!black}İ} TÜBİTAK BiGG kabul/red bilgisi (ortalama 3 ay süre).
\end{itemize}

\subsection*{Ay 4 — Pilot Canlı + ML Baseline}

\begin{itemize}
    \item {\color{blue!50!black}T} İlk friendly pilot ``shadow mode'' (sadece veri toplama).
    \item {\color{blue!50!black}T} XGBoost yük tahmini baseline. Walk-forward validation.
    \item {\color{blue!50!black}T} Algoritma + tahmin entegrasyon: rolling 1-saat tahmin → karar.
    \item {\color{red!50!black}İ} 1. yasal tarama: KVKK aydınlatma metni canlı.
    \item {\color{green!50!black}P} 2 ücretli pilot ön görüşme.
\end{itemize}

\subsection*{Ay 5 — Pilot Optimizasyonu + Pre-Seed Hazırlık}

\begin{itemize}
    \item {\color{blue!50!black}T} Pilot ``control mode'' aç. A/B test.
    \item {\color{blue!50!black}T} İlk ``öncesi/sonrası'' raporu çıkar (1 ay veri).
    \item {\color{red!50!black}İ} Pre-seed pitch deck. Finansal model (3 yıllık).
    \item {\color{green!50!black}P} Galata Business Angels, ScaleX, KEIRETSU İstanbul ile 5--10 toplantı.
\end{itemize}

\subsection*{Ay 6 — İlk Ücretli Müşteri + Pre-Seed Sözleşme}

\begin{itemize}
    \item {\color{green!50!black}P} İlk ücretli pilot sözleşmesi imza (Setup fee + 6 ay SaaS).
    \item {\color{red!50!black}İ} Pre-seed kapanış görüşmeleri --- 3--8M TL hedef.
    \item {\color{blue!50!black}T} Multi-tenant hardening (RLS, tenant izolasyonu).
    \item {\color{blue!50!black}T} OCPP tarafında 3+ farklı charger üretici testi (ABB, Phihong, Schneider).
\end{itemize}

\subsection*{Ay 7--8 — Pre-Seed Kapandı + Ekip Büyüme}

\begin{itemize}
    \item {\color{red!50!black}İ} Pre-seed para hesapta. Şirketleşme (Ltd Şti $\rightarrow$ A.Ş. üzerine düşün).
    \item {\color{red!50!black}İ} Hire: 1 backend dev, 1 frontend dev, 1 sales dev rep (SDR).
    \item {\color{blue!50!black}T} 2--3 ücretli pilot canlı. CS workflow oturmuş.
    \item {\color{blue!50!black}T} V3 sürüm: alarm, rapor otomatik, KVKK ROPA, ISO 27001 başlangıcı.
    \item {\color{green!50!black}P} 5--10 yeni outbound müşteri görüşmesi.
\end{itemize}

\subsection*{Ay 9--10 — Ölçeklendirme + Demand Response}

\begin{itemize}
    \item {\color{blue!50!black}T} OCPP 2.0.1 desteği başla.
    \item {\color{blue!50!black}T} Demand response prototipi --- BEDAŞ pilot için hazırlık.
    \item {\color{blue!50!black}T} LSTM tahmin V3 araştırma. Akademik partner imza.
    \item {\color{green!50!black}P} 5+ ödeyen müşteri.
    \item {\color{red!50!black}İ} TÜBİTAK 1507 KOBİ Ar-Ge başvuru.
\end{itemize}

\subsection*{Ay 11--12 — Seed Hazırlık + Diferansiyatör}

\begin{itemize}
    \item {\color{red!50!black}İ} Seed pitch deck v2. Yatırımcı listesi (212, Re-Pie, Earlybird, Diffusion, Speedinvest).
    \item {\color{blue!50!black}T} PV/ESS proto demo (1 saha PV + sim BESS).
    \item {\color{blue!50!black}T} Harmonic analiz proto demo (Janitza verileri).
    \item {\color{blue!50!black}T} V4 Mimari: multi-region DR, KKTC saha denemesi.
    \item {\color{green!50!black}P} 8--12 ödeyen müşteri, 1--2.5M TL ARR.
\end{itemize}

% ===================================================================
\section{Donanım Gereksinimleri}

\subsection{Lab Kurulumu (1 kez, Faz 1)}

\begin{tabularx}{\textwidth}{l X r}
\toprule
\textbf{Ekipman} & \textbf{Açıklama} & \textbf{Yaklaşık Maliyet (TL)} \\
\midrule
Vestel Bonus DC60 & Test charger (60 kW) & 250.000 \\
veya 2. el ABB Terra HPC & 2nd hand 50 kW & 150.000 \\
Test bataryası / EV emülatör & Charger çıkış yüklemek için & 80.000 (kira: 5K/ay) \\
Janitza UMG 605 & Power Quality Analyzer & 35.000 \\
Endüstriyel switch & MOXA EDS-405A & 5.000 \\
IoT Gateway & Siemens IOT2050 / Raspberry Pi 4 & 4.000 \\
Akım trafosu (CT) seti & Saha çıkışı + her charger & 6.000 \\
Pano malzeme & Sigorta, kablaj, montaj & 10.000 \\
Switching işçilik & Elektrikçi & 15.000 \\
\midrule
\textbf{Toplam (orta tahmin)} & & \textbf{$\sim$310.000} \\
\textbf{Toplam (kira-tabanlı minimum)} & & \textbf{$\sim$80.000 + 5K/ay} \\
\bottomrule
\end{tabularx}

\textbf{Tavsiye:} Vestel Bonus'u sıfır almak yerine \textbf{Vestel ile ``demo birim ödünç''} müzakeresi yap. Pilotta kullan, satış kanalı oluşturursan iade et / al.

\subsection{Pilot Saha Donanımı (saha başına)}

\begin{itemize}
    \item Janitza UMG 96-PA (kompakt sayaç) --- 12K TL.
    \item IoT gateway --- 4K TL.
    \item 4G yedek bağlantı (M2M SIM, Turkcell IoT) --- 200 TL/ay.
    \item Kurulum işçilik --- 15K TL.
    \item Toplam saha kurulum: $\sim$40K TL — pilotta setup fee'den karşılanır.
\end{itemize}

% ===================================================================
\section{Yazılım Stack'i (Önerilen)}

\begin{tabularx}{\textwidth}{l X}
\toprule
\textbf{Katman} & \textbf{Teknoloji} \\
\midrule
Backend & Python 3.12, FastAPI, SQLModel, Alembic, Celery + Redis (queue) \\
OCPP & \texttt{mobilityhouse/ocpp} (Python) — 1.6 + 2.0 destek \\
Algoritma & \texttt{simulation\_v3.py} → \texttt{evlb/algorithms/} paket; numpy, scipy, CVXPY (Faz 3 MILP) \\
ML & scikit-learn, XGBoost, LightGBM (V2). PyTorch + transformers (V3+). MLflow (model registry). \\
DB & PostgreSQL 16 + TimescaleDB extension; Redis 7 \\
Streaming & NATS (hafif) veya Kafka (saha sayısı 10+) \\
Frontend & Next.js 14 + TypeScript + TanStack Query + Tailwind + shadcn/ui \\
Mobile & React Native (faz 3) \\
IaC & Terraform (Hetzner provider) + Ansible \\
CI/CD & GitHub Actions; Docker; ghcr.io \\
Observability & Grafana + Prometheus + Loki; Sentry \\
Auth & Keycloak (self-hosted) veya Auth0 \\
Test & pytest + pytest-asyncio; Playwright (E2E); Locust (load test) \\
Container & Docker Compose (dev), Kubernetes (k3s pilot, EKS prod) \\
\bottomrule
\end{tabularx}

% ===================================================================
\section{Akademik İşbirliği \& Araştırma}

\subsection{Hedef Kurumlar}

\begin{tabularx}{\textwidth}{l X}
\toprule
\textbf{Kurum} & \textbf{Olası işbirliği} \\
\midrule
İTÜ Elektrik Mühendisliği / Enerji Enstitüsü & Yük tahmini, harmonik analizi, MILP \\
ODTÜ EE / EE-MS programı & RL ajan tezi, demand response \\
Boğaziçi Sürdürülebilir Kalkınma & Ekonomik analiz, müşteri davranışı \\
Yıldız EE & Güç kalitesi, harmonik filtre tasarımı \\
TOGG R\&D & V2G pilot, ISO 15118 \\
TÜBİTAK MAM Enerji Enstitüsü & Şebeke entegrasyonu, BESS yönetimi \\
KOÇ EE / Mühendislik & Optimizasyon, makine öğrenmesi \\
Sabancı CDS\&E & ML / forecasting \\
\bottomrule
\end{tabularx}

\subsection{Yapı Önerileri}
\begin{itemize}
    \item Her kurumdan 1 öğretim üyesini ``advisor'' olarak shareholder'la (\%0.25--0.5 shadow equity).
    \item TÜBİTAK 1505 (Üniversite-Sanayi Ortak Ar-Ge) --- 1.5--2.5M TL, ortak proje formatı.
    \item Yüksek lisans tezi sponsor (3K TL/ay × 18 ay = 54K TL) --- ucuz Ar-Ge gücü.
    \item Yıllık IEEE Smart Grid TR konferans katılım + makale.
\end{itemize}

% ===================================================================
\section{Riskler ve Atladığımız Açıklar}

\subsection{Atlanmış 10 Risk Daha}

\begin{enumerate}
    \item \textbf{Charger üreticilerinin OCPP yorumu farklı.} Standart 1.6'da boşluklar var; üretici X'in \texttt{SetChargingProfile} algılayışı üretici Y'den farklı. Lab testi şart, tek üretici tek pilot.
    \item \textbf{Saha elektrikçi ile görüşme gerekli.} Sayaç bağlantısı için AG/OG yetkili elektrikçi şart. KOSGEB elektrikçi listesi.
    \item \textbf{Saha sahibinin elektrikçisi panik yapar.} ``Sayaca müdahale ediyorlar'' algısı. Çözüm: pilot ön-toplantısında elektrikçiyi dahil et.
    \item \textbf{Kurumsal müşteri AŞ değilse vergi/finans kötü.} Faz 2'de A.Ş.'ye geç (avukat: 30K TL maliyet).
    \item \textbf{Para bekleme süresi}. Kurumsal vade 60--90 gün; cash flow problem. \textbf{Setup fee'yi peşin al}, kontrata yaz.
    \item \textbf{Pilot LOI'nin sözleşmeye çevrilememesi.} LOI = niyet, sözleşme = tahsis. Pilot kontratı şablonu hazır olsun.
    \item \textbf{Algoritma yanlış karar verirse}. Sigorta lazım. Aynı zamanda ``pilot evresinde no liability'' kontrat maddesi.
    \item \textbf{Charger arızası bizim hatamız sanılır.} Sözleşmede ``EVSE arızasından sorumlu değiliz'' maddesi. Telemetri ile kanıt logları.
    \item \textbf{Yatırımcı iletişim yorgunluğu.} Pre-seed ile Seed arası 6--9 ay; haberleşme + güncellemeler şart. Aylık investor update e-postası.
    \item \textbf{Ekip içi anlaşmazlık (3 kurucu).} Founders' Agreement gün 1; tahkim maddesi; vesting cliff.
\end{enumerate}

% ===================================================================
\section{30 / 60 / 90 / 180 / 365 Gün Mihenk Taşları}

\begin{tabularx}{\textwidth}{l X}
\toprule
\textbf{Tarih} & \textbf{Olması Gereken} \\
\midrule
30. gün & Founders' Agreement imzalı; 10+ customer discovery; OCPP 1.6 spec okundu; mobilityhouse/ocpp lab kurulumu çalışıyor. \\
60. gün & 30 customer discovery + 3 LOI; TÜBİTAK BiGG başvurusu yapıldı; Vestel Bonus sipariş edilmiş; FastAPI skeleton hazır. \\
90. gün & İlk gerçek charger BootNotification alındı; tarife modülü çalışıyor; ROI calculator demo; 1 friendly pilot saha kurulum başladı. \\
180. gün & 1 friendly + 1--2 ücretli pilot canlı; XGBoost tahmin V2 üretimde; pre-seed kapanış görüşmeleri; 200--500K TL ARR pipeline. \\
365. gün & 8--12 ödeyen müşteri; 1--2.5M TL ARR; 100--300 charger yönetilen; 5--12 kişilik ekip; Seed turu görüşmeleri başlamış. \\
\bottomrule
\end{tabularx}

% ===================================================================
\section{Doğrulama / Verification}

\subsection{Aylık Check-in (her ayın 1'i)}

\begin{itemize}
    \item Customer discovery sayısı / hedef.
    \item LOI sayısı / hedef.
    \item Pilot sayısı (friendly + ücretli) / hedef.
    \item ARR / hedef.
    \item Yönetilen charger sayısı.
    \item Yönetilen toplam kapasite (kW).
    \item Tahmin doğruluğu (MAPE 1h, 24h).
    \item Algoritma ROI (\% tasarruf, $\Delta$bekleme).
    \item Cash runway (ay).
\end{itemize}

\subsection{Yön Değiştirme Tetikleri (Red Flags)}

\begin{itemize}
    \item 6 haftada 30 görüşme tamamlanmadıysa $\rightarrow$ operasyon disiplini sorunu.
    \item 3 ayda 1 ödeyen pilot yoksa $\rightarrow$ ICP yanlış, geri dön.
    \item 6 ayda 100K TL ARR'a ulaşılmadıysa $\rightarrow$ değer önerisi sorgula.
    \item 9 ayda 3 ödeyen müşteri yoksa $\rightarrow$ pivot ya da kapatma sinyali.
    \item Pilotta ``öncesi/sonrası'' raporu \%20 tasarruf altındaysa $\rightarrow$ algoritma değiştir.
    \item Tahmin MAPE 24h $>$ \%25 üstündeyse $\rightarrow$ feature engineering, ICP daraltma.
\end{itemize}

\subsection{Ürün KPI'ları (üretimde)}

\begin{tabularx}{\textwidth}{l l X}
\toprule
\textbf{KPI} & \textbf{Hedef} & \textbf{Ölçüm} \\
\midrule
Trafo aşım dakikası & $-\%80$ vs baseline & power\_log $>$ limit \\
Aylık fatura tasarrufu & $\%20$--$\%40$ & TEDAŞ tarifesi $\times$ delta tüketim \\
Charger başına saatlik servis & $\times 1.5$--$2.5$ & transaction\_count / hour / charger \\
Müşteri başına ortalama bekleme & $-\%30$ & wait\_minutes ortalama \\
Tahmin MAPE 1h / 24h & $<\%10$ / $<\%18$ & walk-forward CV \\
Sistem uptime & $\%99.5$ & Prometheus uptime \\
OCPP bağlantı drop'u & $<\%1$ aylık & connection\_drops / heartbeat\_count \\
\bottomrule
\end{tabularx}

% ===================================================================
\section{Kapanış --- Kritik 5 Şey}

\begin{enumerate}
    \item \textbf{Customer discovery'yi atlama.} 6 hafta disiplinli geçir. Türkiye'de teknik kurucular bu adımı atlar ve 6 ay sonra ``ürün doğru ama müşteri yanlış'' acısı yaşar.

    \item \textbf{OCPP'yi öğren ama overengineering yapma.} 1.6 Core + Smart Charging Profile yeterli; 2.0.1 Faz 3. ISO 15118 Faz 4. Spec'i derinleşmek vakit yer; pratik lab testi daha hızlı öğretir.

    \item \textbf{Tahmin modelinde XGBoost yeterlidir.} V1 baseline → V2 XGBoost → Production. Transformer/LSTM yatırımcı pitch'inde göster ama production'a XGBoost koy. \textbf{Veri olmadan model olmaz}; ilk 3 ay veri toplama.

    \item \textbf{3 kurucu = 3x velocity ama 3x koordinasyon riski.} Founders' Agreement gün 1, vesting 4 yıl 1 yıl cliff, haftalık 1:1 sync, aylık strateji review, paylaşımlı OKR.

    \item \textbf{Hızlı gitmenin tek limiti hukuk + regülatör.} EPDK + KVKK + KİK alanlarını ihmal etme. 2. ayda 1 günlük yasal danışman ile yol haritası temizle (Esin Avukatlık veya Pekin Bayar Türkiye).
\end{enumerate}

\vspace{1cm}
\begin{center}
\textit{Bu doküman \texttt{simulation\_v3.py} prototipinin Türkiye'de ticarileşme yol haritasıdır.\\
Yazıldığı tarih: \today\\
Sürüm: v1 — Birleşik Teknik \& İş Planı, 12 Ay, 3 Kurucu.}
\end{center}

\end{document}
